% --- TEMPLATE TRABAJO SOCIAL (UNESUM) ---
% Estándar de Reporteo Académico Limpio y Desacoplado
\documentclass[12pt,spanish]{article}

\usepackage{graphicx}
\usepackage{geometry}
\usepackage{indentfirst} 
\usepackage{caption}
\usepackage{floatrow} 

% Forzar estilos de caption antes de cualquier carga
\floatsetup[figure]{capposition=top}
\floatsetup[table]{capposition=top}

\usepackage{booktabs}
\usepackage{float}
\usepackage{longtable}
\usepackage{threeparttable}
\usepackage{array}
\usepackage{calc}
\usepackage{fontspec}
\setmainfont{Times New Roman}
\usepackage{amsmath, amssymb}
\usepackage{fancyhdr}
\usepackage{setspace} 
\usepackage{ragged2e} 
\usepackage{titlesec}
\usepackage{url}
\def\UrlBreaks{\do\/\do-\do\.\do\_\do\?\do\=\do\&\do\%}
\usepackage[colorlinks,linkcolor=blue,urlcolor=blue,citecolor=blue,breaklinks]{hyperref}

% --- Pandoc Compatibility Block ---
\providecommand{\tightlist}{%
  \setlength{\itemsep}{0pt}\setlength{\parskip}{0pt}}

% --- Pandoc Citeproc (CSL) compatibility ---
\newlength{\cslhangindent}
\setlength{\cslhangindent}{1.27cm} 
\newlength{\csllabelwidth}
\setlength{\csllabelwidth}{3em}
\newcommand{\citeproctext}{} 
\newcommand{\citeprocitem}[2]{\item} 
\newenvironment{CSLReferences}[2] 
 {\begin{list}{}{
  \setlength{\leftmargin}{\cslhangindent}
  \setlength{\parsep}{0pt}
  \setlength{\itemsep}{6pt}
  \ifodd #1
   \setlength{\leftmargin}{\leftmargin + \csllabelwidth}
   \setlength{\itemindent}{-\csllabelwidth}
  \fi
  }
  \renewcommand{\bibitem}[2][]{\item}
 }
 {\end{list}}
\newcommand{\CSLBlock}[1]{#1\hfill\break}
\newcommand{\CSLLeftMargin}[1]{\parbox[t]{\csllabelwidth}{#1}}
\newcommand{\CSLRightInline}[1]{\parbox[t]{\linewidth - \csllabelwidth}{#1}\break}
\newcommand{\CSLIndent}[1]{\hspace{\cslhangindent}#1}

\makeatletter
\@ifundefined{paragraph}{\newcounter{paragraph}}{}
\@ifundefined{subparagraph}{\newcounter{subparagraph}}{}
\makeatother

% Solución al error 'No counter none defined' en longtable
\makeatletter
\newcounter{none}
\@ifpackageloaded{caption}{
  \DeclareCaptionType{none}
  \captionsetup[none]{name=Tabla}
}{}
\makeatother

\AtBeginDocument{
  \renewcommand{\tablename}{Tabla}
  \renewcommand{\figurename}{Figura}
  \renewcommand{\refname}{Referencias}
  \justifying
  \setlength{\parindent}{1.27cm}
  \setlength{\parskip}{0pt}
}

% Márgenes APA 7
\geometry{
  a4paper,
  margin=25.4mm,
  headheight=15mm
}

% --- Encabezados ---
\pagestyle{fancy}
\fancyhf{} 
\renewcommand{\headrulewidth}{0pt}
\fancyfoot[R]{\thepage}

\fancypagestyle{plain}{
  \fancyhf{}
  \fancyfoot[R]{\thepage}
  \renewcommand{\headrulewidth}{0pt}
}

% --- Títulos ---
\titleformat{\section}{\normalfont\normalsize\bfseries\centering}{\thesection}{1em}{}
\titleformat{\subsection}{\normalfont\normalsize\bfseries}{\thesubsection}{1em}{}
\titleformat{\subsubsection}{\normalfont\normalsize\bfseries\itshape}{\thesubsubsection}{1em}{}
\titleformat{\paragraph}{\normalfont\normalsize\bfseries}{\theparagraph}{1em}{}
\titlespacing*{\section}{0pt}{12pt}{6pt}
\titlespacing*{\subsection}{0pt}{12pt}{6pt}
\titlespacing*{\subsubsection}{0pt}{12pt}{6pt}
\titlespacing*{\paragraph}{0pt}{12pt}{6pt}

\captionsetup[table]{labelfont=bf, textfont=bf, font=small}\makeatletter\@ifpackageloaded{caption}{}{\usepackage{caption}}\AtBeginDocument{%
\ifdefined\contentsname
  \renewcommand*\contentsname{Table of contents}
\else
  \newcommand\contentsname{Table of contents}
\fi
\ifdefined\listfigurename
  \renewcommand*\listfigurename{List of Figures}
\else
  \newcommand\listfigurename{List of Figures}
\fi
\ifdefined\listtablename
  \renewcommand*\listtablename{List of Tables}
\else
  \newcommand\listtablename{List of Tables}
\fi
\ifdefined\figurename
  \renewcommand*\figurename{Figure}
\else
  \newcommand\figurename{Figure}
\fi
\ifdefined\tablename
  \renewcommand*\tablename{Table}
\else
  \newcommand\tablename{Table}
\fi
}\@ifpackageloaded{float}{}{\usepackage{float}}
\floatstyle{ruled}
\@ifundefined{c@chapter}{\newfloat{codelisting}{h}{lop}}{\newfloat{codelisting}{h}{lop}[chapter]}
\floatname{codelisting}{Listing}\newcommand*\listoflistings{\listof{codelisting}{List of Listings}}\makeatother\makeatletter\makeatother\makeatletter\@ifpackageloaded{caption}{}{\usepackage{caption}}
\@ifpackageloaded{subcaption}{}{\usepackage{subcaption}}\makeatother

\begin{document}

% --- CARÁTULA UNESUM ---
\begin{center}
    \vspace*{1.5cm}
    \normalsize\textbf{UNIVERSIDAD ESTATAL DEL SUR DE MANABÍ} \\[0.2cm]
    \normalsize Creada mediante la Ley Pública Registro Oficial 261 del 7 de febrero del 2001 \\[0.4cm]
    \normalsize\textbf{FACULTAD DE CIENCIAS SOCIALES, HUMANÍSTICAS Y DE LA EDUCACIÓN} \\[0.3cm]
    \textbf{CARRERA TRABAJO SOCIAL} \\[1.2cm]
    
    \textbf{NOMBRE DE LA MATERIA:} Fundamentos del Trabajo
Social \\[0.3cm]
    \textbf{SEMESTRE Y PARALELO:} Primer Semestre ``C'' \\[0.3cm]
    \textbf{PERIODO ACADÉMICO:} PI 2026 \\[1.2cm]
    
    \textbf{TEMA:} \\[0.2cm]
    Análisis sobre el Funcionalismo, Marxismo y Teorías
Interpretativas \\[1.2cm]
    
    \textbf{INTEGRANTE:} \\[0.2cm]
    Sandra Lorena Seraquive Cuenca \\[1.2cm]
    
    \textbf{DOCENTE FACILITADOR:} \\[0.2cm]
    Lcda. Ligia Geomara Párraga Vélez, Mgs. \\[1.8cm]
    
    \textbf{JIPIJAPA, MANABÍ, ECUADOR}
\end{center}

\thispagestyle{empty}
\newpage
\pagenumbering{arabic}
\doublespacing
\setlength{\parindent}{1.27cm}

\section{Introducción}\label{introducciuxf3n}

El Trabajo Social utiliza las teorías sociológicas para orientar la
intervención. A continuación, se examinarán las limitaciones operativas
y conceptuales del funcionalismo, el marxismo y las teorías
interpretativas, orden en el cual serán analizados para identificar sus
alcances prácticos.

\section{Desarrollo}\label{desarrollo}

Bajo la perspectiva funcionalista, la principal limitación del Trabajo
Social radica en la despolitización de las demandas sociales, al
priorizar la contención burocrática y la adaptación del sujeto sobre el
cambio estructural (Lomelí Gutiérrez et al. 2023). Este enfoque reduce
la práctica a una intervención instrumental que fragmenta los problemas
colectivos en expedientes individuales. Como contraste a este enfoque
integrador, la Table~\ref{tbl-comparativa} detalla y contrapone las
dimensiones de estos marcos teóricos con el modelo crítico e
interpretativo.

\begingroup
\singlespacing
\fontsize{9pt}{11pt}\selectfont
\renewcommand{\arraystretch}{0.9}
\setlength{\tabcolsep}{3pt}

\begin{longtable}[]{@{}
  >{\raggedright\arraybackslash}p{(\linewidth - 6\tabcolsep) * \real{0.2500}}
  >{\raggedright\arraybackslash}p{(\linewidth - 6\tabcolsep) * \real{0.2500}}
  >{\raggedright\arraybackslash}p{(\linewidth - 6\tabcolsep) * \real{0.2500}}
  >{\raggedright\arraybackslash}p{(\linewidth - 6\tabcolsep) * \real{0.2500}}@{}}
\caption{Comparación de enfoques
teóricos}\label{tbl-comparativa}\tabularnewline
\toprule\noalign{}
\begin{minipage}[b]{\linewidth}\raggedright
\textbf{Dimensión}
\end{minipage} & \begin{minipage}[b]{\linewidth}\raggedright
\textbf{El Funcionalismo y la Adaptación Institucional}
\end{minipage} & \begin{minipage}[b]{\linewidth}\raggedright
\textbf{El Marxismo y la Intervención Crítica}
\end{minipage} & \begin{minipage}[b]{\linewidth}\raggedright
\textbf{Las Teorías Interpretativas y la Construcción de Sentido}
\end{minipage} \\
\midrule\noalign{}
\endfirsthead
\toprule\noalign{}
\begin{minipage}[b]{\linewidth}\raggedright
\textbf{Dimensión}
\end{minipage} & \begin{minipage}[b]{\linewidth}\raggedright
\textbf{El Funcionalismo y la Adaptación Institucional}
\end{minipage} & \begin{minipage}[b]{\linewidth}\raggedright
\textbf{El Marxismo y la Intervención Crítica}
\end{minipage} & \begin{minipage}[b]{\linewidth}\raggedright
\textbf{Las Teorías Interpretativas y la Construcción de Sentido}
\end{minipage} \\
\midrule\noalign{}
\endhead
\bottomrule\noalign{}
\endlastfoot
\textbf{Concepción de la sociedad} & Sistema integrado de partes
interdependientes (Alonso Benito 2024). & Estructura conflictiva
dividida por clases (Tejerina Montaña 2024). & Construcción
intersubjetiva cotidiana (Flores Solano 2023). \\
\textbf{Origen del problema social} & Desajustes respecto a las normas
vigentes. & Desigualdades materiales y opresión de clase. & Ruptura o
conflicto en la atribución de sentido diario. \\
\textbf{Rol y objetivo profesional} & Gestor de recursos orientado al
equilibrio y ajuste del sujeto (Lomelí Gutiérrez et al. 2023). &
Promotor de organización colectiva para transformar estructuras de
poder. & Acompañante de procesos reflexivos para la autogestión del
bienestar (Juan Toset and Soraya Mustafa 2023). \\
\end{longtable}

\endgroup

Por su parte, el enfoque crítico marxista enfrenta el riesgo del
determinismo macroestructural (Tejerina Montaña 2024). Esta postura
puede postergar la respuesta operativa inmediata ante la urgencia
cotidiana del sujeto, al supeditar la práctica profesional
exclusivamente a la transformación revolucionaria del sistema
capitalista (Juan Toset and Soraya Mustafa 2023).

Por otro lado, las teorías interpretativas centran su análisis en la
subjetividad cotidiana. No obstante, su limitación radica en el
aislamiento microsocial (Flores Solano 2023), omitiendo la coerción de
las estructuras de poder y las condiciones materiales objetivas que
condicionan la autonomía individual.

\section{Conclusiones}\label{conclusiones}

La tensión entre adaptación (funcionalismo), transformación (marxismo) y
significación (teorías interpretativas) exige una praxis reflexiva. El
Trabajo Social debe intervenir críticamente en los márgenes
institucionales mediante la gestión de recursos viables, sin descuidar
el cuestionamiento a las desigualdades estructurales. Así, la
intervención parte de la narrativa cotidiana de los sujetos, incidiendo
en lo inmediato sin renunciar al cambio estructural.

\section*{Referencias}\label{referencias}
\addcontentsline{toc}{section}{Referencias}

\protect\phantomsection\label{refs}
\begin{CSLReferences}{1}{1}
\bibitem[\citeproctext]{ref-alonso2024funcionalismo}
Alonso Benito, Luis Enrique. 2024. {``Funcionalismo y Paradigmas
Sociológicos. Anotaciones Críticas a Un Libro de Jacques
Coenen-Huther.''} \emph{Revista Española de Investigaciones
Sociológicas}, no. 37: 207--17.
\url{https://doi.org/10.5477/cis/reis.37.207}.

\bibitem[\citeproctext]{ref-flores2023diversidad}
Flores Solano, José Alberto. 2023. {``Diversidad y Tejido Social: Una
Aproximación Desde El Interaccionismo Simbólico.''} \emph{CONfines de
Relaciones Internacionales y Ciencia Política} 19 (36): 24--48.
\url{https://doi.org/10.46530/cf.vi36/cnfns.n36.p24-48}.

\bibitem[\citeproctext]{ref-juantoset2023enfoque}
Juan Toset, Elena María, and Pamela Soraya Mustafa. 2023. {``El Enfoque
Crítico Del Trabajo Social Bajo La Perspectiva de La Teoría Social de
Marx: Una Propuesta Al Trabajo Social.''} \emph{Servicios Sociales y
Política Social} 40 (129): 27--38.
\url{https://www.serviciossocialesypoliticasocial.com/-120}.

\bibitem[\citeproctext]{ref-lomeli2023formacion}
Lomelí Gutiérrez, Ruth, Victoria Colmenares Ríos, and Noemí Macedonio
Toledo. 2023. {``Formación de Los Profesionales Del Trabajo Social y La
Intervención Con Familia: De Lo Funcional Estructural a La Propuesta
Socio-Constructivista.''} \emph{Universidad Veracruzana}, 219--31.
\url{https://www.uv.mx}.

\bibitem[\citeproctext]{ref-tejerina2024conflicto}
Tejerina Montaña, Benjamín. 2024. {``Las Teorías Sociológicas Del
Conflicto Social. Algunas Dimensiones Analíticas a Partir de k. Marx y
g. Simmel.''} \emph{Revista Española de Investigaciones Sociológicas},
no. 55: 47--63. \url{https://doi.org/10.5477/cis/reis.55.47}.

\end{CSLReferences}

\end{document}
